\documentclass[11pt]{article}
\usepackage[T1]{fontenc}
\usepackage{lmodern,amsmath,amssymb,amsthm}
\usepackage[margin=1in]{geometry}
\usepackage{hyperref}
\hypersetup{hidelinks,hypertexnames=false,pdftitle={EGNO Theorem 3.6.4: Universal grading},pdfauthor={Pavel Etingof; Shlomo Gelaki; Dmitri Nikshych; Victor Ostrik}}
\newcommand{\Z}{\mathbb{Z}}
\newcommand{\Zp}{\mathbb{Z}_+}
\newcommand{\Aad}{A_{\mathrm{ad}}}
\newcommand{\sourceloc}[1]{\par\smallskip\noindent\textit{Source: #1.}\par\nopagebreak[4]\smallskip}
\newcommand{\editorial}[1]{\par\smallskip\noindent\textsf{Editorial note.} #1\par\smallskip}
\title{The Adjoint Based Subring and Universal Grading\\Theorem 3.6.4}
\author{Pavel Etingof \and Shlomo Gelaki \and Dmitri Nikshych \and Victor Ostrik}
\date{Private reading transcription of the 2015 author final copy\\with a separately identified 2026 author clarification}
\begin{document}
\maketitle

% BEGIN CONTEXT EGNO:3.6.4
\noindent\textbf{Source and selected scope.}
\emph{Tensor Categories}, AMS Mathematical Surveys and Monographs 205
(2015), author final copy, 362 physical pages. The sole primary selection
is Theorem 3.6.4 on physical page 77, printed page 60, with its complete
printed proof on physical page 78, printed page 61. Propositions 3.6.2,
Lemma 3.6.3, and Corollary 3.6.6 are complete uncounted support and
consequences of the same construction. The original opening unital
context, the later wording ``based ring,'' and the finite-rank aside
are retained in their original positions. No global finite-rank
condition is inserted.

\noindent The source is
\url{https://math.mit.edu/~etingof/egnobookfinal.pdf}.
The separate author corrections dated February 14, 2026 are retained at
\url{https://math.mit.edu/~etingof/tensbookcor.pdf}.
The selected clarification of the based-module dependency is displayed
separately below. It is not silently substituted into the 2015 text.

\subsection*{Selected definitions and coefficient identity}
\sourceloc{Physical page 66, printed page 49; Definitions 3.1.1 and 3.1.3, Remark 3.1.2 and equation (3.1)}
Let $\Zp$ denote the semi-ring of non-negative integers.

\noindent\textbf{Definition 3.1.1.}
Let $A$ be a ring which is free as a $\Z$-module.
\begin{enumerate}
\renewcommand{\labelenumi}{(\roman{enumi})}
\item A $\Zp$-basis of $A$ is a basis $B=\{b_i\}_{i\in I}$ such that
$b_i b_j=\sum_{k\in I}c^k_{ij}b_k$, where $c^k_{ij}\in\Zp$.
\item A $\Zp$-ring is a ring with a fixed $\Zp$-basis and with identity $1$
which is a non-negative linear combination of the basis elements.
\item A unital $\Zp$-ring is a $\Zp$-ring such that $1$ is a basis element.
\end{enumerate}

\noindent\textbf{Remark 3.1.2.}
Every $\Zp$-ring is assumed to have an identity, but according to our
terminology a $\Zp$-ring is not necessarily unital.

Let $A$ be a $\Zp$-ring, and let $I_0$ be the set of $i\in I$ such that
$b_i$ occurs in the decomposition of $1$. Let $\tau:A\to\Z$ denote the
group homomorphism defined by
\begin{equation}
\tau(b_i)=\begin{cases}1&\text{if }i\in I_0,\\0&\text{if }i\notin I_0.\end{cases}
\tag{3.1}
\end{equation}

\noindent\textbf{Definition 3.1.3.}
A $\Zp$-ring $A$ with basis $\{b_i\}_{i\in I}$ is called a based ring if
there exists an involution $i\mapsto i^*$ of $I$ such that the induced map
\[
a=\sum_{i\in I}a_i b_i\longmapsto a^*=\sum_{i\in I}a_i b_{i^*},
\qquad a_i\in\Z,
\]
is an anti-involution of the ring $A$, and
\begin{equation}
\tau(b_i b_j)=\begin{cases}1&\text{if }i=j^*,\\0&\text{if }i\ne j^*.\end{cases}
\tag{3.2}
\end{equation}

\sourceloc{Physical pages 66--67, printed pages 49--50; complete Proposition 3.1.6 and its printed proof, followed by Definition 3.1.7}
\noindent\textbf{Proposition 3.1.6.}
\emph{In any based ring, the number $c^{k^*}_{ij}$ is invariant under
cyclic permutations of $i,j,k$.}
\begin{proof}
It is easy to see that $c^{k^*}_{ij}=\tau(b_i b_j b_k)$, and
$\tau(xy)=\tau(yx)$. Thus, $c^{k^*}_{ij}$ is cyclically symmetric.
\end{proof}

\noindent\textbf{Definition 3.1.7.}
A multifusion ring is a based ring of finite rank. A fusion ring is a
unital based ring of finite rank.

\sourceloc{Physical page 73, printed page 56; complete Definitions 3.4.1 and 3.4.2}
\noindent\textbf{Definition 3.4.1.}
Let $A$ be a $\Zp$-ring with basis $\{b_i\}$. A $\Zp$-module over $A$ is
an $A$-module $M$ with a fixed $\Z$-basis $\{m_l\}_{l\in L}$ such that all
the structure constants $a^k_{il}$ (defined by the equality
$b_i m_l=\sum_k a^k_{il}m_k$) are non-negative integers.

The direct sum of $\Zp$-modules is also a $\Zp$-module whose basis is the
union of the bases of the summands. We say that a $\Zp$-module is
indecomposable if it is not isomorphic to a nontrivial direct sum of
$\Zp$-modules.

\noindent\textbf{Definition 3.4.2.}
A $\Zp$-module $M$ over a $\Zp$-ring $A$ is called irreducible if it has
no proper $\Zp$-submodules (in other words, the $\Z$-span of any proper
subset of the basis of $M$ is not an $A$-submodule).

\sourceloc{Physical page 74, printed page 57; original dependency, Exercise 3.4.3(i)}
\noindent\textbf{Exercise 3.4.3(i), as printed in 2015.}
Prove that a $\Zp$-module over a based ring is irreducible if and only
if it is indecomposable.

\sourceloc{Separate corrections dated February 14, 2026, physical page 2; selected Section 3.4 clarification and first correction to Exercise 3.4.3(i)}
\begin{quote}
\textit{Section 3.4.} The definition of a based $\Zp$-module over a based
ring is accidentally omitted, although it is used further in the text
(first time in Proposition 3.6.2).

By definition, a based $\Zp$-module over a based ring $A$ is one where
if $x\in A$ acts by a matrix $X$ then $x^*$ acts by $X^T$.

\textit{Exercise 3.4.3(i).} The $\Zp$-module should assumed be based.
\end{quote}
\editorial{This is the selected based-module clarification relevant to
the later based-module text. The full separately retained corrections
also contain the authors' non-based examples and an alternative
finite-rank condition with Lemma 2.1 and its proof on page 3. That
alternative is not imposed as a global finite-rank hypothesis here.}

\sourceloc{Physical page 75, printed page 58; complete Definition 3.5.1 and the following grading/faithfulness paragraphs, before the fusion-ring dimension specialization}
Let $A$ be a unital based ring with a $\Zp$-basis $B=\{b_i\}_{i\in I}$.

\noindent\textbf{Definition 3.5.1.}
Let $G$ be a group. A grading of $A$ by $G$ is a partition
$B=\bigsqcup_{g\in G}B_g$ of the basis of $A$ into disjoint subsets such
that for all $g,h\in G$ and for all basis elements $b_i\in B_g$ and
$b_j\in B_h$ the product $b_i b_j$ is a $\Zp$-linear combination of
elements of $B_{gh}$.

Note that if $b_i\in B_g$ then $b_i^*\in B_{g^{-1}}$. We will call $B_g$
the component of $B$ corresponding to $g\in G$, and $B_1$ (where $1$
denotes the identity element of $G$) the trivial component of $B$.
Note that the trivial component includes all basis elements contained
in the decomposition of $b_i b_i^*$, $i\in I$. In particular, the trivial
component contains $b_i$ for all $i\in I_0$.

If $A$ is $G$-graded, we have a decomposition
\[
A=\bigoplus_{g\in G}A_g,
\]
where $A_g$ is the $\Z$-span of $B_g$, $g\in G$. We will say that $A_g$ is
the component of $A$ corresponding to $g$. The subring $A_1$ will be
called the trivial component of $A$.

The grading is called faithful if all the components $B_g$, $g\in G$,
are nonempty.
% END CONTEXT EGNO:3.6.4

\section*{3.6. The adjoint based subring and universal grading}
\sourceloc{Physical page 77, printed page 60; original section opening and all support before Theorem 3.6.4}
% BEGIN CONSTRUCTION SETUP EGNO:3.6.4
Let $A$ be a unital based ring with a $\Zp$-basis $B=\{b_i\}_{i\in I}$.

\noindent\textbf{Definition 3.6.1.}
The adjoint subring $\Aad\subset A$ is the minimal based subring of $A$
with the property that $b_i b_i^*$ belongs to $\Aad$ for all $i\in I$;
i.e., $\Aad$ is generated by all basis elements of $A$ contained in
$b_i b_i^*$, $i\in I$.

If $A$ is of finite rank (i.e., a fusion ring), let us define
$J:=\sum_{i\in I}b_i b_i^*$ (note that this is a central element of $A$
by Proposition 3.1.8 for $x=1$). Then $\Aad$ can be alternatively defined
as the $\Z$-linear span of basis elements contained in $J^n$,
$n=0,1,2,\ldots$.
% END CONSTRUCTION SETUP EGNO:3.6.4

% BEGIN SUPPORT STATEMENT EGNO:3.6.2
\noindent\textbf{Proposition 3.6.2.}
\emph{Any one-sided (i.e., left or right) based $\Aad$-submodule
$M\subset A$ of a based ring $A$ is automatically an $\Aad$-sub-bimodule.}
% END SUPPORT STATEMENT EGNO:3.6.2
% BEGIN SUPPORT PROOF EGNO:3.6.2
\begin{proof}
It suffices to prove the statement for left modules. We have
\[
b_k b_i b_{i^*}
 =\sum_{r:b_r\in b_k b_i}c^r_{ki}b_r b_{i^*}
 =\sum_{r:b_r\in b_k b_i}c^{i^*}_{r^*k}b_r b_{i^*}
\]
(where $b_r\in b_k b_i$ means that $b_r$ occurs in the decomposition of
$b_k b_i$). On the other hand,
\[
\sum_{r,i}c^{i^*}_{r^*k}b_r b_{i^*}
 =\left(\sum_r b_r b_{r^*}\right)b_k.
\]
But
\[
\sum_{r,i}c^{i^*}_{r^*k}b_r b_{i^*}
 =\sum_{r:b_r\in b_k b_i}c^{i^*}_{r^*k}b_r b_{i^*}+E_{ki},
\]
where $E_{ki}$ is a non-negative linear combination of basis elements.
Hence,
\[
\left(\sum_r b_r b_{r^*}\right)b_k=b_k b_i b_{i^*}+E_{ki}.
\]
So if $b_k\in M$ then $b_k b_i b_{i^*}\in M$ for all $i$, as desired.
\end{proof}
% END SUPPORT PROOF EGNO:3.6.2

% BEGIN DECOMPOSITION CONTEXT EGNO:3.6.4
Let $A$ be a based ring. We can view $A$ as a based $\Zp$-bimodule over
$\Aad$. As such, it decomposes into a direct sum of indecomposable
based $\Aad$-bimodules: $A=\bigoplus_{a\in G}A_a$, where $G$ is the
index set. This decomposition is unique up to a permutation of $G$.
We may assume that there is an element $1\in G$ such that $A_1=\Aad$.
Note that $(A_a)^*=\{X^*\mid X\in A_a\}$, $a\in G$, is an
indecomposable based $\Aad$-submodule of $A$ and hence
$(A_a)^*=A_{a^*}$ for some $a^*\in G$.
% END DECOMPOSITION CONTEXT EGNO:3.6.4

% BEGIN SUPPORT STATEMENT EGNO:3.6.3
\noindent\textbf{Lemma 3.6.3.}
\emph{For all $x,y\in A_a$, $a\in G$, we have $xy^*\in\Aad$.}
% END SUPPORT STATEMENT EGNO:3.6.3
% BEGIN SUPPORT PROOF EGNO:3.6.3
\begin{proof}
We may assume that $x,y$ are basis vectors. Observe that
\[
\begin{split}
M_y=\operatorname{span}\{X\in A_a\mid{}&X\text{ is contained in }
b_{i_1}b_{i_1}^*\cdots b_{i_n}b_{i_n}^*y\\
&\text{for some }n,i_1,\ldots,i_n\}
\end{split}
\]
is a based left $\Aad$-submodule of $A_a$. Since $A_a$ is an
indecomposable based $\Aad$-bimodule, by Exercise 3.4.3(i), it is
irreducible, and so it follows from Proposition 3.6.2 that it is in
fact irreducible as a based left $\Aad$-module. Thus, $M_y=A_a$.
Therefore, $x\in M_y$, so $x$ is contained in
$b_{i_1}b_{i_1}^*\cdots b_{i_n}b_{i_n}^*y$ for some
$n,i_1,\ldots,i_n$, i.e., $xy^*$ is a summand in
$b_{i_1}b_{i_1}^*\cdots b_{i_n}b_{i_n}^*yy^*\in\Aad$, as required.
\end{proof}
% END SUPPORT PROOF EGNO:3.6.3

\sourceloc{Physical page 77, printed page 60; complete Theorem 3.6.4}
% BEGIN PRIMARY STATEMENT EGNO:3.6.4
\noindent\textbf{Theorem 3.6.4.}
\emph{There is a canonical group structure on the index set $G$ with
the multiplication defined by the following property:}
\begin{equation}
\begin{gathered}
ab=c\quad\text{if and only if }x_a x_b\in A_c,\\
\text{for all }x_a\in A_a,\ x_b\in A_b,\ a,b,c\in G.
\end{gathered}
\tag{3.6}\label{eq:3.6}
\end{equation}
\emph{The identity of $G$ is $1$ and the inverse of $a\in G$ is $a^*$.}
% END PRIMARY STATEMENT EGNO:3.6.4

\sourceloc{Physical page 78, printed page 61; complete printed proof}
% BEGIN PRIMARY PROOF EGNO:3.6.4
\begin{proof}
We need to check that the binary operation in \eqref{eq:3.6} is well defined.
Let $a,b\in G$ and let $x_a,y_a\in A_a$, $x_b,y_b\in A_b$ be basis
elements of $A$.

Suppose that the product $x_a x_b$ contains a basis element $x_c\in A_c$
and the product $y_a y_b$ contains a basis element $y_d\in A_d$ for
$c\ne d$. By Lemma 3.6.3, there is a positive integer $n$ and indices
$i_1,\ldots,i_n$ such that the element
$z:=x_a b_{i_1}b_{i_1}^*\cdots b_{i_n}b_{i_n}^*y_a^*$ contains
$x_a x_b y_b^*y_a^*=(x_a x_b)(y_a y_b)^*$.
But by Proposition 3.6.2, the left based $\Aad$-module generated by
$x_a$ is also a right based $\Aad$-module, which means that
$x_a b_{i_1}b_{i_1}^*\cdots b_{i_n}b_{i_n}^*$ is a summand in a sum of
elements of the form $b_{j_1}b_{j_1}^*\cdots b_{j_m}b_{j_m}^*x_a$,
and hence $z\in A_1$. Therefore, $Y:=x_c y_d^*$ is in $A_1$.
Multiplying both sides of the last equality by $y_d$ on the right, we
conclude that $Yy_d$ is in $A_c\cap A_d$, a contradiction.

Thus, $x_a x_b$ and $y_a y_b$ both belong to the same component $A_c$
and so the binary operation \eqref{eq:3.6} is well-defined. It is easy
to see that it defines a group structure on $G$.
\end{proof}
% END PRIMARY PROOF EGNO:3.6.4

\sourceloc{Physical page 78, printed page 61; complete Definition 3.6.5 and Corollary 3.6.6, stopping before the upper-central-series paragraph}
% BEGIN UNIVERSAL PROPERTY CONTEXT EGNO:3.6.4
\noindent\textbf{Definition 3.6.5.}
We will call the grading $A=\bigoplus_{a\in G}A_a$ constructed in
Theorem 3.6.4 the universal grading of $A$. The group $G$ will be
called the universal grading group of $A$ and denoted by $U(A)$.
% END UNIVERSAL PROPERTY CONTEXT EGNO:3.6.4

% BEGIN SUPPORT STATEMENT EGNO:3.6.6
\noindent\textbf{Corollary 3.6.6.}
\emph{Every based ring $A$ has a canonical faithful grading by the
group $U(A)$. Any other faithful grading of $A$ by a group $G$ is
determined by a surjective group homomorphism $\pi:U(A)\to G$.}
% END SUPPORT STATEMENT EGNO:3.6.6
% BEGIN SUPPORT PROOF EGNO:3.6.6
\begin{proof}
Let $A=\bigoplus_{g\in G}A^g$ be a faithful grading of $A$.
Since for every basis element $X\in A$ we have $XX^*\in A^1$, it
follows that $A^1$ contains $\Aad$ as a based subring.
Hence, each $A^g$ is a based $\Aad$-submodule of $A$.
This means that every component $A_a$, $a\in U(A)$, of the universal
grading $A=\bigoplus_{a\in U(A)}A_a$ of $A$ belongs to some
$A^{\pi(a)}$ for some well-defined $\pi(a)\in G$.
Clearly, the map $a\mapsto\pi(a)$ is a surjective homomorphism.
\end{proof}
% END SUPPORT PROOF EGNO:3.6.6

% BEGIN CREDIT CONTEXT EGNO:3.6.4
\sourceloc{Physical page 81, printed page 64; selected complete bibliographical note 3.6}
The notion of the adjoint subring of a fusion ring is motivated by the
notion of the adjoint subcategory of a fusion category, introduced in
[ENO2, 8.5], and developed in [GelaN]. The material on the adjoint subring,
the universal grading, and nilpotent based rings is taken from [GelaN].

\sourceloc{Physical page 347, selected original bibliography entry [GelaN]}
\noindent[GelaN] S. Gelaki and D. Nikshych. \emph{Nilpotent fusion categories}.
Advances in Mathematics \textbf{217} (2008), no. 3, 1053--1071.

\editorial{The four book authors provide the collected exposition;
the selected note preserves its explicit historical source credit.
The complete printed proofs above keep their original compressed
observations and final group-axiom check. The initial unital setup
and subsequent general wording remain visible together, as do the
2015 dependency and the separate 2026 author clarification.}
% END CREDIT CONTEXT EGNO:3.6.4
\end{document}
